% Source PDF: source/普通高考/2013/2013湖北文.pdf, page 1, question 13.
% PDF truth: mutool/pdfimages review; no embedded bitmap, vector flowchart.
% Node→directed-edge list:
%   开始→输入m→(A=1,B=1,i=0)→(i=i+1)→(A=A×m)→(B=B×i)→(A<B?)；
%   否→回到 i=i+1 前的主干连接点；是→输出i→结束。
% All outlines/connectors are solid. Labels and formulas are centered in local zero-indent nodes.

\begin{tikzpicture}[
  >=Stealth,
  x=1cm,
  y=1cm,
  line width=0.45pt,
  line cap=round,
  line join=round,
  font=\normalsize,
  flowtext/.style={anchor=center, align=center,
    execute at begin node={\setlength{\parindent}{0pt}}},
  flowbox/.style={flowtext, draw, minimum height=0.68cm,
    inner xsep=4pt, inner ysep=2pt},
  terminal/.style={flowbox, rounded corners=2pt, minimum width=1.35cm},
  io/.style={flowbox, trapezium, trapezium left angle=72, trapezium right angle=108},
  process/.style={flowbox, minimum width=2.45cm},
  decision/.style={flowbox, diamond, aspect=2.7, inner sep=0pt,
    minimum width=4.55cm, minimum height=1.10cm},
  branchlabel/.style={flowtext, inner sep=0pt},
]
  \node[terminal] (start) at (0,10.80) {开始};
  \node[io, minimum width=2.00cm] (input) at (0,9.40) {输入 $m$};
  \node[flowbox, minimum width=3.95cm] (init) at (0,8.00) {$A=1,B=1,i=0$};
  \node[process, minimum width=2.05cm] (inc) at (0,6.62) {$i=i+1$};
  \node[process, minimum width=2.55cm] (amul) at (0,5.28) {$A=A\times m$};
  \node[process, minimum width=2.55cm] (bmul) at (0,3.94) {$B=B\times i$};
  \node[decision] (test) at (0,2.55) {$A<B?$};
  \node[io, minimum width=1.75cm] (output) at (0,1.15) {输出 $i$};
  \node[terminal] (finish) at (0,-0.15) {结束};

  \draw[->] (start.south) -- (input.north);
  \draw[->] (input.south) -- (init.north);
  \draw[->] (init.south) -- (inc.north);
  \draw[->] (inc.south) -- (amul.north);
  \draw[->] (amul.south) -- (bmul.north);
  \draw[->] (bmul.south) -- (test.north);

  % The 否 branch is an independent right-side loop and returns leftward
  % into the main trunk above i=i+1; the vertical trunk arrow remains separate.
  \draw[line width=0.45pt] (test.east) -- (3.25,2.55) -- (3.25,7.20);
  \draw[->] (3.25,7.20) -- (0,7.20);
  \node[branchlabel] at (2.72,2.92) {否};

  \draw[->] (test.south) -- (output.north);
  \node[branchlabel] at (0.32,1.96) {是};
  \draw[->] (output.south) -- (finish.north);
  \path[use as bounding box] (-0.65,-0.55) rectangle (3.55,11.15);
\end{tikzpicture}
